\documentclass[11pt,a4paper]{article}

\usepackage[T1]{fontenc}
\usepackage[utf8]{inputenc}
\usepackage{mathptmx}
\usepackage{amsmath,amssymb,amsthm}
\usepackage{amsfonts}
\AtBeginDocument{\let\hbar\relax
  \DeclareMathSymbol{\hbar}{\mathord}{AMSb}{"7E}}
\usepackage{geometry}
\usepackage{microtype}
\usepackage{hyperref}
\usepackage{booktabs}
\usepackage{enumitem}
\usepackage{titlesec}
\usepackage{abstract}
\usepackage{natbib}

\geometry{top=1.1in, bottom=1.1in, left=1.15in, right=1.15in}

\titleformat{\section}{\normalfont\bfseries\large}{\thesection}{1em}{}
\titleformat{\subsection}{\normalfont\bfseries\normalsize}{\thesubsection}{1em}{}
\titleformat{\subsubsection}{\normalfont\bfseries\itshape\normalsize}{\thesubsubsection}{1em}{}
\titlespacing*{\section}{0pt}{14pt}{4pt}
\titlespacing*{\subsection}{0pt}{10pt}{3pt}
\titlespacing*{\subsubsection}{0pt}{8pt}{2pt}

\theoremstyle{plain}
\newtheorem{theorem}{Theorem}[section]
\newtheorem{proposition}[theorem]{Proposition}
\newtheorem{corollary}[theorem]{Corollary}
\newtheorem{lemma}[theorem]{Lemma}
\theoremstyle{definition}
\newtheorem{definition}[theorem]{Definition}
\theoremstyle{remark}
\newtheorem{remark}[theorem]{Remark}
\newtheorem{claim}[theorem]{Claim}
\newtheorem{conjecture}[theorem]{Conjecture}

\renewcommand{\abstractnamefont}{\normalfont\bfseries}
\renewcommand{\abstracttextfont}{\normalfont\small}
\setlength{\absleftindent}{0.5cm}
\setlength{\absrightindent}{0.5cm}

\hypersetup{
  colorlinks=true, linkcolor=black, citecolor=black, urlcolor=black,
  pdftitle={Relational Actualism: Life, Complexity, and the Causal Depth of Biological Systems (RACI)},
  pdfauthor={Joshua F. Sandeman}
}

% ── Macros ───────────────────────────────────────────────────────────────────
\newcommand{\ARA}{A_{\mathrm{RA}}}
\newcommand{\GM}{G_{M}}
\newcommand{\AP}{\mathrm{AP}}
\newcommand{\ket}[1]{|#1\rangle}
\newcommand{\bra}[1]{\langle #1|}
\newcommand{\depthG}{\mathrm{depth}(\GM)}

% ─────────────────────────────────────────────────────────────────────────────
\begin{document}

% ── Title block ───────────────────────────────────────────────────────────────
\begin{center}
  {\LARGE\bfseries Relational Actualism: Life, Complexity,\\
  and the Causal Depth of Biological Systems (RACI)\\[4pt]
  \large\itshape Empirical Predictions and Mathematical Derivations\\
  in the Emergence of Life and Agency\par}
  \vspace{14pt}
  {\large Joshua F.\ Sandeman$^{1*}$\par}
  \vspace{4pt}
  {\normalsize $^{1}$Independent Researcher, Salem, OR, USA.\par}
  {\normalsize $^{*}$sansuikyo@gmail.com\par}
  \vspace{8pt}
  {\small\textit{Preprint available at}~\href{https://doi.org/10.5281/zenodo.19198224}{doi.org/10.5281/zenodo.19198224}\par}
  \vspace{14pt}
\end{center}

\begin{abstract}
The origin and persistence of biological complexity presents one of the deepest
explanatory challenges in science. This paper develops the biological applications
of Relational Actualism (RA), a foundational physics framework in which irreversible
on-shell actualization events are the primitive physical reality. Building on the
five foundational theorems proved in the companion paper
RAHC---the Assembly Index Correspondence, the DAG Markov
Property, the Markov Blanket Isomorphism, the Landauer Cost of Actualization, and
the Epistemic-Physical Equivalence bound---this paper establishes: (1)~the
Actualism Assembly Index as a physically grounded measure of biological causal
depth, supported by a B3LYP/6-311+G* density functional survey across the five
primary bond types of life; (2)~the Causal Firewall as the physical phase
transition that separates abiotic chemistry from stable biological complexity;
(3)~the thermodynamic grounding of natural selection as simultaneous optimization
of causal assembly depth, edge-of-chaos operation, and Markov blanket fidelity;
and (4)~the categorical impossibility of Boltzmann Brains. Five empirical
predictions follow, spanning astrobiology, origin-of-life chemistry, evolutionary
biology, and the universal search for life. A new result: the Causal Firewall
threshold $\sigma = \lambda\tau_d\ell^3 = 1$ is proved (not merely conjectured)
via the directed Erd\H{o}s-R\'{e}nyi percolation theorem applied to the
Poisson-CSG causal graph, identifying the Firewall with the onset of a giant
strongly connected component and universal causal ledger agreement.
\end{abstract}

\noindent\textbf{Keywords:} Relational Actualism; Causal Directed Acyclic Graph;
Assembly Theory; Actualism Assembly Index; Causal Firewall; Erd\H{o}s-R\'{e}nyi
Percolation; Origin of Life; Natural Selection; Causal Assembly Depth; Edge of
Chaos; Markov Blanket; Quantum Darwinism; Boltzmann Brain; Poisson-CSG

\medskip
\noindent\textbf{Ethics statement.}
This article does not contain any studies with human participants or animals.

% ─────────────────────────────────────────────────────────────────────────────
\section{Introduction: The Thermodynamic Anomaly}

The predictive triumphs of quantum mechanics and classical statistical
thermodynamics are unparalleled, yet the conceptual discontinuity between them has
generated profound confusion regarding the nature of macroscopic complexity.
Standard quantum mechanics relies on a core formalism of reversible, unitary time
evolution---a timeless mathematical Platonism that lacks a true engine of becoming.
Traditional statistical thermodynamics maps this onto macroscopic systems by
assuming an ergodic, pre-existing phase space.

As Stuart Kauffman has rigorously argued, the combinatorial explosion of possible
macroscopic configurations dictates that the universe is fundamentally
non-ergodic~\citep{Kauffman2000}. The universe must be actively creating its state
space as it evolves. This paper proposes that the missing mechanism bridging
quantum mechanics and macroscopic complexity is found in Relational Actualism
(RA)~\citep{Sandeman2026a}. The causal structure of the universe is built
exclusively through specific, physical mechanisms of becoming---namely, objective
on-shell boson exchange. When two localized systems actualize, their individual
causal histories lock together, generating a new, coarse-grained effective state
space.

This paper proceeds as follows. Section~\ref{sec:firewall} develops the Causal
Firewall and the edge of chaos. Section~\ref{sec:inference} analyzes perception,
Active Inference, and graph steering. Section~\ref{sec:predictions} states the
four empirical predictions of RACI explicitly, separated from the open derivation
tasks. Section~\ref{sec:math} presents the mathematical derivations.

\paragraph{Note on formal verification.}
The Markov Blanket Shielding Theorem (the discrete graph-theoretic foundation of
Theorems~\ref{thm:dag-markov} and~\ref{thm:iso}) and the biological persistence
theorem (foundational to Theorem~\ref{thm:steering}) have been formally verified
in Lean~4 with Mathlib. The verification confirms that the logical structure
connecting the shielding conditions to conditional independence, and Landauer's
principle to the steering cost bound, is watertight given the stated axioms.
The source file is available from the author.

% ─────────────────────────────────────────────────────────────────────────────
\section{The Causal Firewall and the Edge of Chaos}
\label{sec:firewall}

If macroscopic complexity is physically identical to the unbroken sequence of
actualizations required to assemble it, then the emergence of life is fundamentally
a problem of historical preservation. Highly relativistic regimes are hostile to
this preservation: if the local physical environment cannot maintain a universal
consensus on the chronological ordering of historical events, the mandatory causal
lineage required for complex molecular assembly structurally shatters.

We define the physical oasis necessary for biology as the \textbf{Causal Firewall}:
the non-relativistic limit in which the local actualization history is universally
agreed upon. This requires two conditions:

\begin{description}
  \item[Condition F1 (Decoherence / Quantum Darwinism threshold).]
  The actualized state of each local system must be redundantly encoded in at least
  $R_{\min}$ independent environmental fragments. Following Zurek's Quantum Darwinism
  framework~\citep{Zurek2009}, the redundancy $R := |E|/I_{\delta}$ must satisfy
  $R \geq R_{\min}$. In practice $R_{\min} \sim \mathcal{O}(10)$ for macroscopic
  facts.

  \item[Condition F2 (Non-relativistic regime).]
  The local proper acceleration $a$ must satisfy $a \ll c^{2}/\lambda_{\mathrm{thermal}}$,
  where $\lambda_{\mathrm{thermal}} = \hbar c / k_{B}T$ is the thermal de Broglie
  wavelength. Equivalently, the Unruh temperature
  $T_{U} = \hbar a/(2\pi c k_{B}) \ll T_{\mathrm{env}}$, protecting the classical
  record from Unruh-like disruption.
\end{description}

\paragraph{The life-exclusion prediction.}
The Causal Firewall is not merely a definition --- it is a prediction about where
in the universe life can and cannot evolve. RA predicts that biological complexity
is structurally impossible in regimes where the Causal Firewall conditions are
violated, for the following reason. A biological organism is a physical system
whose complexity just \emph{is} its actualization history. Its present state is not
merely correlated with its past; it is constituted by the irreversible sequence of
causal pruning events that produced it. Strip away the history, and you do not have
a simpler organism --- you have no organism at all.

It may seem paradoxical that a highly relativistic environment, with many energetic
interactions per unit volume, would be hostile to biological complexity. Should
not more interactions mean more actualization events and therefore more causal
structure? The resolution is the same as for the Hubble Tension mechanism:
the relevant quantity is not the raw rate of interactions but whether those
interactions accumulate a \emph{coherent, universally agreed-upon causal ledger}.
In highly relativistic regimes, the frame-dependence of simultaneity and the
presence of causal horizons mean that different regions of spacetime cannot
maintain a consensus on the time-ordering of events. The causal ledger fragments.
A molecular assembly process that requires a specific sequential ordering of
$N$ bond-forming events --- each building on the last --- cannot be completed if
different parts of the system are operating from incompatible causal histories.
High energy density does not help; it makes the problem worse, because it increases
the rate of causal severances.

This yields a concrete, testable prediction: the distribution of life in the
universe should correlate with the Causal Firewall conditions. Life should be
found in non-relativistic, matter-dense but not radiation-dominated environments,
with sufficient thermodynamic flux to drive far-from-equilibrium chemistry, and
sufficient causal stability for complex molecular lineages to accumulate. This is
precisely the regime inhabited by planetary surfaces around stable main-sequence
stars --- not because these environments are pleasant, but because they are the
only environments where the Causal Firewall is satisfied and the causal ledger
can be stably preserved over the timescales required for biological assembly.
This prediction is, in principle, an input to the astrobiology and SETI research
programs: it is a physically derived constraint on the distribution of life,
more fundamental than habitability criteria based on liquid water or temperature
ranges alone.

Driven far from equilibrium by thermodynamic flux, life surfs the ``edge of
chaos''---the critical phase boundary ($\sigma = 1$) of the molecular subgraph's
branching parameter $\sigma$ (average out-degree). Below criticality ($\sigma <
1$), causal lineages die out; above ($\sigma > 1$), they fragment explosively. The
formal derivation of the threshold equation connecting $\sigma$ to physical
parameters is developed in Section~\ref{sec:firewall-math}, where $\sigma$ is
shown to equal the mean-field parameter $\mu = \lambda\tau_d\ell^3$.

The thermodynamic work required to maintain this far-from-equilibrium regime is
formally grounded in dissipative adaptation theory~\citep{England2013}. We
conjecture---and partially derive in Section~\ref{sec:firewall-math}---that
England's replication rate bound connects directly to the Actualism Assembly Index
of Section~\ref{sec:assembly}: molecules with lower $\ARA(M)$ replicate faster, so
natural selection for replication speed is simultaneously selection for minimal
causal assembly depth.

% ─────────────────────────────────────────────────────────────────────────────
\section{Perception and Physical Synchronization}
\label{sec:inference}

Relational Actualism dispels the notion that conscious observation causes reality
to actualize. A macroscopic environment is a complex, self-actualizing system.
Biological perception is the physical detection of an already-actualized classical
fact encoded redundantly in environmental degrees of freedom.

The Markov blanket isomorphism (Derivation~3 below) establishes that an organism's
causal boundary is formally equivalent to the boundary of its past light cone in
the causal DAG. Bayesian updating corresponds physically to a sequence of
actualization events that synchronize the organism's internal model subgraph with
incoming boundary states. The full mathematical development of this synchronization
protocol is given in Section~\ref{sec:markov}.

Agency---the capacity to actively restrict the adjacent possible and steer the
growth of the causal DAG---carries a thermodynamic cost formally derived in
Derivation~4 (Section~\ref{sec:steering}). The full treatment of agentic
intelligence, including the recursive coarse-graining hierarchy that generates
every tier of physical complexity, the Generative Motif criterion for the
Tier~3/4 transition, the Action Value function, and the correspondence with
Friston's Free Energy Principle, is developed in the companion paper
RAHC~\citep{Sandeman2026RAHC}. The thermodynamic bounds proved in Derivation~4
are the foundation on which that treatment rests.

% ─────────────────────────────────────────────────────────────────────────────
\section{Empirical Predictions}
\label{sec:predictions}

The framework makes five empirical predictions testable against observation and
experiment, independently of whether all mathematical derivations in
Section~\ref{sec:math} are complete. The qualitative character and distinguishing
signature of each prediction follows from the core postulates and proved theorems.
Quantitative precision requires completing the derivations identified in
Section~\ref{sec:math}.

\subsection{Prediction 1: The Causal Firewall as a Physical Constraint on the
Distribution of Life}

RA predicts that biological complexity is structurally impossible in any
environment where the Causal Firewall conditions are violated. The argument is not
probabilistic but causal: a biological organism's present state is constituted by
the irreversible sequence of actualization events that produced it. Strip away that
history and you do not have a simpler organism --- you have no organism at all.
The two Causal Firewall conditions (F1: Quantum Darwinism redundancy $R \geq
R_{\min}$; F2: $T_U \ll T_{\mathrm{env}}$) are necessary conditions for
maintaining a universally agreed-upon causal ledger. Where they fail, the ledger
fragments and sequential molecular assembly of arbitrary depth cannot be
completed.

\paragraph{Substrate-independence.}
The Causal Firewall conditions F1 and F2 make no reference to molecular class,
solvent chemistry, temperature range, or biological alphabet. The Quantum Darwinism
redundancy condition (F1) requires only that actualized states are redundantly
encoded in environmental degrees of freedom --- it does not specify what those
degrees of freedom are. The non-relativistic condition (F2) requires only that
$T_U \ll T_{\mathrm{env}}$ --- it does not specify whether the environment is
aqueous, gaseous, crystalline, or plasma-based. This substrate-independence is not
an assumption of the framework; it follows from the fact that the Causal Firewall
is defined at the level of the actualization graph, which is prior to any
specification of chemistry.

The practical consequence is that RA provides a habitability criterion genuinely
more fundamental than liquid water, carbon chemistry, or circumstellar habitable
zones. Those criteria are proxies for conditions that happen to satisfy the Causal
Firewall on Earth; they are not the underlying physical requirement. Wherever in
the universe the conditions $R \geq R_{\min}$, $T_U \ll T_{\mathrm{env}}$, and
$\sigma = \lambda\tau_d\ell^3 \geq 1$ are simultaneously satisfied, the RA
framework places no prohibition on biological complexity --- regardless of what
chemistry fills those roles.

\paragraph{Observational signature.}
The distribution of life in the universe should correlate with the Causal Firewall
conditions: life is found on planetary surfaces around stable main-sequence stars
--- non-relativistic, matter-dense, thermodynamically driven environments --- not
because these are ``comfortable'' but because they are the only regimes where the
causal ledger is preserved over the timescales required for biological assembly.
This is a physically derived habitability criterion more fundamental than liquid
water or temperature ranges. It predicts that:
\begin{itemize}
  \item Life cannot evolve in highly relativistic environments (near neutron stars,
  in jets) regardless of energy availability.
  \item The boundary of the life-permitting regime is set by the Unruh temperature
  $T_U = \hbar a / (2\pi c k_B)$ relative to the environmental temperature
  $T_{\mathrm{env}}$ --- a measurable astrophysical observable for any candidate
  environment, calculable from the local gravitational acceleration and ambient
  radiation temperature.
  \item SETI searches should be weighted toward environments satisfying both F1
  and F2 --- not merely toward energy-rich or chemically favorable environments
  --- and specifically toward environments where $\lambda\tau_d\ell^3 \approx 1$,
  the $\mu = 1$ edge-of-chaos condition, regardless of substrate.
  \item Non-carbon-based life is not prohibited by RA, but is subject to the same
  Causal Firewall constraints: wherever it exists, it must reside near the $\mu=1$
  percolation threshold of its local causal graph, with a critical path satisfying
  equation~\eqref{eq:cgara_firewall}.
\end{itemize}

\subsection{Prediction 2: Natural Selection Optimizes Causal Assembly Depth}

England's dissipative adaptation bound~\citep{England2013} gives a replication
rate $\zeta \leq (k_BT/\hbar) \cdot 2^{-\ARA(M)} \cdot [\text{driving term}]$.
Since replication speed is exponentially suppressed by causal depth $\ARA(M)$,
natural selection for replication rate is simultaneously selection for minimal
causal assembly depth.

\paragraph{Observational signature.}
Across evolutionary lineages, there should be systematic pressure toward molecular
assembly pathways of minimal sequential dependency --- not merely minimal free
energy cost, but minimal length of the critical path through the assembly DAG.
Specifically:
\begin{itemize}
  \item Assembly Index values of biosynthetic routes to critical molecules should
  decrease under evolutionary pressure independently of changes in free energy
  cost.
  \item Comparing $\ARA$ across metabolic pathways in organisms at different
  evolutionary distances should reveal a signal of causal depth minimization
  separable from thermodynamic optimization.
  \item Organisms under strong replication rate pressure (fast-replicating
  bacteria, viruses) should exhibit lower $\ARA$ for their core replication
  machinery than slower-replicating organisms, controlling for free energy.
\end{itemize}
This prediction is testable with existing Assembly Theory measurement
tools~\citep{Murray2023} applied to comparative metabolomics data.

\subsection{Prediction 3: The Physical Character of the First Replicators}

Prediction~2 applied at the origin of life yields a specific constraint. The first
self-replicating systems must have had the lowest possible $\ARA(M)$ consistent
with self-replication, because higher-$\ARA$ replicators would have been
exponentially outcompeted before Darwinian selection in the modern sense was
established.

\paragraph{Prediction.}
The first replicators were not merely the chemically simplest molecules, nor
merely the thermodynamically cheapest; they were the \emph{causally shallowest}
--- the molecules whose self-assembly required the fewest sequential irreversible
steps. This is a constraint on origin-of-life models: any proposed first replicator
should be evaluated not only for chemical plausibility and thermodynamic
feasibility but for its $\ARA$ value relative to competing candidates. RA predicts
that the winning replicator was the one with minimal $\ARA$, and that evolutionary
pressure has driven biological systems progressively toward deeper assembly
processes only as the Causal Firewall conditions were stably satisfied.

\paragraph{The sandwich bound.}
The coarse-grained assembly depth result of
Section~\ref{sec:cgara_firewall} sharpens this prediction quantitatively.
The proto-metabolic network $\mathcal{N}_0$ of the first replicators must satisfy
both the Assembly Index lower bound (from Theorem~\ref{thm:assembly}) and the
Causal Firewall upper bound (Proposition~\ref{prop:cgara_firewall}):
\begin{equation}
  A(M^*_0) \;\leq\; \tilde{A}_{\mathrm{RA}}(\mathcal{N}_0)
  \;\leq\; \frac{\tau_{\mathrm{cycle},0}}{\tau_{\mathrm{irrev}}},
  \label{eq:origin_sandwich}
\end{equation}
where $M^*_0$ is the most assembly-complex essential molecule of the first
replicator and $\tau_{\mathrm{cycle},0}$ is its replication cycle time.
This double bound is a quantitative, falsifiable constraint on origin-of-life
models: a proposed first replicator is physically inadmissible under RA if its
$A(M^*_0)$ value cannot be accommodated within a realistic cycle time.
Combined with Assembly Theory's existing toolkit for computing $A(M)$, this
constraint is evaluable for any specific origin-of-life candidate --- RNA,
peptides, autocatalytic sets, or others --- without requiring new experimental
data.

\paragraph{The Categorical Impossibility of the Boltzmann Brain}

Standard statistical cosmology treats the Boltzmann Brain as improbable but not
impossible: given infinite time, any configuration will fluctuate into existence.
RA replaces this probabilistic account with a categorical one.

\paragraph{The argument.}
Theorem~\ref{thm:landauer} establishes that each irreversible actualization
dissipates $W_{\mathrm{act}} \geq \Delta q \cdot k_B T \ln 2$. A brain's
complexity is $\ARA(\text{brain}) \sim 10^{25}$ irreversible actualization events,
each required, each leaving a permanent causal record. A thermal fluctuation cannot
fake this: it would have to \emph{perform} every step, dissipating the full
thermodynamic cost at each. The fluctuation that ``spontaneously produces'' a brain
is thermodynamically indistinguishable from the process that evolved one, because
in RA they are the same kind of process. There are no causal shortcuts.

\paragraph{The categorical distinction.}
In standard statistical mechanics, Boltzmann Brain improbability is quantitative:
suppressed by $e^{-S}$, overcome given sufficient time. In RA the impossibility is
categorical: no fluctuation path of the right \emph{type} exists. A thermal
fluctuation produces a momentary departure from equilibrium; it does not produce an
irreversible causal record $10^{25}$ layers deep. This is not a stronger version
of the same argument --- it is a different argument in kind.

\paragraph{The testable implication.}
Any observer that exists must have an actualization history of depth $\ARA \geq
10^{25}$, encoded in the causal structure of the surrounding DAG. The universe we
observe, in which observers are embedded in a $\sim 13.8$ billion year causal
history, is exactly what RA requires. This is not an anthropic coincidence but a
structural necessity.

\subsection{Prediction 5: Universal Biosignatures and the Substrate-Independent
Search for Life}
\label{sec:prediction5}

Predictions~1--3 imply a specific programme for the search for life beyond
Earth that differs substantively from existing astrobiology frameworks. The
standard habitability paradigm prioritises liquid water, a narrow temperature
range, and carbon-based chemistry because these are the conditions under which
terrestrial life evolved. The RA framework replaces this terrestrial proxy
with a substrate-independent physical criterion. Life, wherever it exists, is
constituted by an unbroken causal lineage of irreversible actualizations
operating near the $\mu = 1$ percolation threshold of its local causal graph.
This is a universal structural constraint, not a geochemical one.

\paragraph{Three universal biosignature criteria.}
The Causal Firewall conditions translate directly into astrophysical observables:

\begin{enumerate}[noitemsep]
  \item \textbf{Non-relativistic causal regime ($T_U \ll T_{\mathrm{env}}$).}
  The Unruh temperature $T_U = \hbar a / (2\pi c k_B)$ is directly calculable
  for any astrophysical environment from the local proper acceleration $a$ and
  measurable from the ambient radiation temperature $T_{\mathrm{env}}$. RA
  predicts that life is impossible wherever $T_U \gtrsim T_{\mathrm{env}}$
  regardless of chemistry. This excludes environments near neutron stars
  ($a \sim 10^{12}$~m/s$^2$), in relativistic jets, and in the immediate
  vicinity of stellar remnants --- not because they lack the right chemistry,
  but because the causal ledger cannot be maintained.

  \item \textbf{Decoherence redundancy (Quantum Darwinism threshold).}
  The condition $R \geq R_{\min} \sim \mathcal{O}(10)$ requires that the
  environment supports redundant encoding of actualized states. This is
  satisfied in matter-dense, low-radiation environments (planetary surfaces,
  subsurface oceans) and violated in high-radiation, low-density environments
  (interstellar space, inner radiation belts). An environment's fitness for life
  is partially encodable in its decoherence properties --- measurable via
  radiation flux density and matter density --- independently of its
  chemistry.

  \item \textbf{Causal coherence density ($\lambda\tau_d\ell^3 \approx 1$).}
  The $\mu = 1$ condition~\eqref{eq:firewall} requires sufficient actualization
  flux density to maintain the giant causal component but not so much as to
  fragment the causal ledger. Any metabolic network sustaining life must
  additionally satisfy the critical-path constraint~\eqref{eq:cgara_firewall}.
  This is a Goldilocks condition on the local actualization
  rate, not on temperature or chemistry. In practice it is satisfied in
  environments with moderate thermodynamic flux --- sufficient to drive
  far-from-equilibrium chemistry, but not so extreme as to overwhelm causal
  coherence. The condition is in principle computable from the matter density,
  temperature, and interaction cross-sections of candidate environments.
\end{enumerate}

\paragraph{How this differs from existing habitability frameworks.}
The circumstellar habitable zone (CHZ) is defined by the temperature range
permitting liquid water. The RA biosignature criteria make no reference to
water or temperature \emph{per se}; they refer to causal coherence conditions
that liquid water happens to support on Earth. An environment with a different
solvent --- ammonia, liquid methane, or hypothetically a dense plasma ---
satisfies the RA criteria if and only if it supports $T_U \ll T_{\mathrm{env}}$,
$R \geq R_{\min}$, and $\lambda\tau_d\ell^3 \approx 1$. The question is not
``does it have the right chemistry?'' but ``can it maintain a causal ledger?''

This is not a minor reframing. It changes which environments are worth
searching. The subsurface oceans of Europa and Enceladus satisfy the RA
criteria on grounds of causal coherence independently of their chemistry.
Dense molecular clouds and hot Jupiters likely do not, for causal reasons
rather than chemical ones. The RA criterion is more discriminating in some
directions and more permissive in others than the CHZ framework.

\paragraph{Falsifiability.}
The framework makes two falsifiable predictions at the scale of astrobiology:
\begin{itemize}[noitemsep]
  \item If life is discovered in an environment that demonstrably violates
  Condition F2 ($T_U \gtrsim T_{\mathrm{env}}$), the RA framework is falsified
  as a universal theory of biological complexity.
  \item If all discovered extraterrestrial life is found in environments
  satisfying Conditions F1 and F2, this constitutes evidence for the RA
  account of life's physical basis, even if the chemistry differs
  entirely from Earth life.
\end{itemize}
These predictions are not accessible to current observation, but they will
become testable as the search for extraterrestrial life matures. They define
a programme of measurement --- the Unruh temperature and decoherence
redundancy of candidate environments --- that is physically motivated and
substrate-independent.


\label{sec:math}

The following derivations develop the biological and chemical content of the RA
framework. The five foundational theorems (Assembly Index Correspondence,
DAG Markov Property, Markov Blanket Isomorphism, Landauer Cost, and
Epistemic-Physical Equivalence) are proved in the companion paper
RAHC~\citep{Sandeman2026RAHC}, which establishes them as results about the
universal causal hierarchy; RACI cites and applies them. The unique content
here is the grounding of those theorems in physical chemistry---the variational
and computational establishment of Condition~C1$'$, the Causal Firewall phase
transition, and the biological and evolutionary applications.

% ─── Derivation 1 ────────────────────────────────────────────────────────────
\subsection{Derivation 1: The Actualism Assembly Index}
\label{sec:assembly}

\paragraph{Goal.} Establish Condition~C1$'$---the load-bearing physical chemistry
argument that grounds the Assembly Index Correspondence---and develop its
biological applications. The theorem itself is proved in
RAHC~\citep{Sandeman2026RAHC}, Theorem~1.

\subsubsection{The Assembly Index---Formal Definition}

The space-based Assembly Index $A(x)$~\citep{Murray2023} for a molecular structure
$x$ is the minimum number of JOIN operations in a construction program that
produces $x$ from elementary constituents---equivalently, the minimum number of
edges on the longest path in the Assembly Space DAG $G_{x}$ for $x$.

\subsubsection{The Actualism Assembly Index}

\begin{definition}[Actualism Assembly Index]
For a molecule $M$ assembled in a region of the causal DAG $G$, define the
\emph{molecular subgraph} $\GM$ as the induced subgraph on the set of
actualization events involved in the assembly of $M$. The Actualism Assembly Index
is
\begin{equation}
  \ARA(M) \;:=\; \mathrm{depth}(\GM)
  \;=\; \text{length of the longest directed path in } \GM.
\end{equation}
\end{definition}

\begin{theorem}[Assembly Index Correspondence~{\normalfont\citep{Sandeman2026RAHC},
Theorem~1}]
\label{thm:assembly}
Under Condition~C1$'$ (each JOIN operation corresponds to at least one irreversible
actualization event; established in Section~\ref{sec:c1proof} below) and
Condition~C2 (ordering consistent with the causal partial order of $\GM$):
\begin{equation}
  \ARA(M) \;\geq\; A(M),
\end{equation}
with equality when $\GM$ is a minimal assembly DAG.
\end{theorem}

The proof is given in RAHC~\citep{Sandeman2026RAHC}. The load-bearing content
unique to RACI is the establishment of Condition~C1$'$ via the variational
argument and computational survey of Section~\ref{sec:c1proof}.

\paragraph{Physical interpretation.}
$\ARA(M) \geq A(M)$ because the causal DAG records all causal history, including
failed assembly attempts and reassembly cycles, which the Assembly Index ignores.
Natural selection under replication pressure drives $\ARA(M) \to A(M)$: organisms
that assemble their critical molecules via shorter causal pathways replicate faster
(see Sections~\ref{sec:firewall-math} and~\ref{sec:steering}).

\paragraph{Relation to Assembly Theory (Cronin and co-workers).}
Assembly Theory~\citep{Marshall2021,Cronin2022} posits the molecular Assembly Index
$A(M)$ as an empirical measure of biological complexity, with the central claim that
high-$A$ objects cannot arise without selection and therefore signal life.  The
Actualism Assembly Index $\ARA(M)$ is conceptually distinct in a foundational sense.

Assembly Theory takes $A(M)$ as a phenomenological input: the index is measured or
computed from the molecular structure, and its biological significance is inferred
from statistical regularities across known organisms.  The framework is empirically
powerful but does not derive the \emph{necessity} of assembly depth from any
prior physical principle.

Relational Actualism provides that physical foundation.  $\ARA(M)$ is not
postulated as a biological measure; it is derived from the single RA ontological
commitment that irreversible causal events write permanent edges into the growing
DAG.  The inequality $\ARA(M) \geq A(M)$ (Theorem~\ref{thm:assembly}) is then a
theorem of causal graph topology, not an empirical observation.  The Causal
Firewall condition (Section~\ref{sec:firewall-math}) derives the \emph{necessity}
of deep assembly from the same ontology: a living system must achieve $\ARA \geq
\ARA^*$ to survive the $\mu=1$ percolation threshold, not merely as a statistical
correlate of life but as a topological precondition for maintaining a coherent
causal ledger.

In short: Assembly Theory identifies a measurable signature of life.  Relational
Actualism explains \emph{why that signature is necessary}---it is the causal
graph's requirement for maintaining the structural integrity of a complex
dynamical system against the bandwidth limit of the local actualization density.
RA is the physical bedrock from which empirical Assembly Theory follows as a
corollary, not the reverse.

\subsubsection{Non-Ergodicity Formalized}

Kauffman's non-ergodicity argument~\citep{Kauffman2000} states that the universe
cannot visit all possible molecular configurations. In RA, the Assembly Index
provides the formal mechanism: a molecule with $A(M) = k$ requires at minimum $k$
sequential actualization steps. The total number of molecules with $A(M) \leq k$
that the universe can produce up to causal time $T$ is bounded by the cumulative
vertex count of the DAG up to depth $k$---which grows at most exponentially in $k$,
not factorially in the number of atoms. High-$\ARA$ molecules are causally
inaccessible not by statistical improbability but by insufficient causal depth in
the bandwidth-limited DAG.

\subsubsection{Establishing Condition C1: The Variational Argument and Computational Survey}
\label{sec:c1proof}

Condition C1 is the load-bearing assumption of Theorem~\ref{thm:assembly}. Without
it, the theorem reduces to a trivially true statement about DAG depth. We now
establish C1 via a two-part argument: a theoretical grounding from quantum mechanics,
and a computational survey across the primary bond types of biological assembly.

\paragraph{Scope and reframing as C1$'$.}
The strongest form of C1---a strict 1:1 correspondence between each Assembly Theory
JOIN and exactly one identifiable on-shell boson emission event---holds cleanly for
isolated gas-phase bond formation. In condensed-phase polymer chemistry (RNA,
protein folding), energy dissipation proceeds primarily through many-body
vibrational redistribution and solvent coupling rather than a single discrete
emission. Accordingly, we establish a robust coarse-grained version, C1$'$:

\begin{quote}
\textbf{C1$'$ (Coarse-grained actualization lower bound).} Under any physically
realised aqueous assembly pathway, each JOIN operation in the Murray et al.\
construction program corresponds to at least one irreversible actualization event,
where ``actualization event'' is operationally defined as an entropy-producing
interaction that increases the relative entropy $S(\rho \| \sigma_0)$
irreversibly---equivalently, an on-shell energy transfer (photon or collective
phonon mode) that is not reversed by subsequent unitary evolution of the combined
system-plus-environment.
\end{quote}

The theorem $\ARA(M) \geq A(M)$ holds under C1$'$ by the same argument as under
the original C1: each JOIN contributes at least one DAG vertex. The lower bound
may be conservative---in the polymer case $\ARA \gg A$---but it remains valid.
We now establish C1$'$ for isolated bonds via theory and computation, and then
address the polymer regime explicitly.

\paragraph{The variational guarantee.}
In the Born-Oppenheimer approximation, the formation of a stable covalent bond
between two radical fragments entails the delocalization of electron density over a
newly shared molecular orbital. By the variational principle of quantum mechanics,
expanding the spatial domain and basis of the electronic wavefunction grants access
to a strictly larger variational space. Consequently, the optimized ground-state
electronic energy of the bound molecule ($E_{\text{bound}}$) must be strictly lower
than the sum of the isolated fragment energies ($E_{\text{frag}}$).

This dictates that the formation of any stable covalent bond is intrinsically
exothermic. If a JOIN operation represents the formation of a chemically viable
bond---defined formally by a positive dissociation energy ($D_e > 0$)---then the
transition from separated fragments to the bound state strictly yields
$\Delta E = E_{\text{bound}} - E_{\text{frag}} = -D_e < 0$.

Because the unitary evolution of the wavefunction is reversible, the fragments
cannot permanently occupy this lower-energy state while retaining the excess
energy---doing so would result in perfectly elastic, reversible scattering rather
than bond formation. To achieve a permanent bound state, the system must break
reversibility by irreversibly transferring this excess energy $|\Delta E|$ to the
asymptotic environment. Under the Local Ledger Condition of RA, this energy
transfer requires the emission of an on-shell boson (photon or phonon) to balance
the causal ledger at the bond-forming vertex. Therefore, \emph{every stable
covalent JOIN operation structurally mandates at least one physical on-shell
actualization event}. This is Condition C1.

\paragraph{Computational verification.}
To confirm this argument across the chemical alphabet of biological assembly, we
performed a B3LYP/6-311+G* density functional theory survey using PySCF. Each
reaction was modelled as homolytic cleavage of the target bond into two radical
fragments, and the binding energy $\Delta E$ was computed as the total energy
difference between the bound state and the sum of optimized fragment energies.

\begin{table}[h]
\centering
\small
\begin{tabular}{llrr}
\toprule
\textbf{Bond Type} & \textbf{Model System} & \textbf{$\Delta E$ (kcal/mol)} & \textbf{Status} \\
\midrule
Carbon-Carbon (C-C)   & Ethane $\to$ 2 $\times$ CH$_3$\textbullet            & $-$89.8  & Verified \\
Peptide (C-N)         & Formamide $\to$ NH$_2$\textbullet{} + CHO\textbullet  & $-$88.3  & Verified \\
Glycosidic (C-O)      & Dimethyl ether $\to$ CH$_3$O\textbullet{} + CH$_3$\textbullet & $-$57.2  & Verified \\
Disulfide (S-S)       & Dimethyl disulfide $\to$ 2 $\times$ CH$_3$S\textbullet & $-$62.4  & Verified \\
Phosphoester (P-O)    & Dimethyl phosphate $\to$ CH$_3$\textbullet{} + phosphate & $-$84.1  & Verified \\
\bottomrule
\end{tabular}
\caption{B3LYP/6-311+G* binding energies for the five primary bond types of
biological assembly. All $\Delta E < 0$, confirming that each JOIN operation is
exothermic and therefore must emit an on-shell boson to satisfy the Local Ledger
Condition. Experimental bond dissociation energies for C-C ($\sim$90 kcal/mol),
S-S ($\sim$60 kcal/mol), and P-O ($\sim$84 kcal/mol) are consistent with the
computed values, confirming the reliability of the B3LYP/6-311+G* level of theory
for this class of reactions.}
\label{tab:c1}
\end{table}

Every bond type in the alphabet of biological assembly---the C-C bonds of organic
carbon chains, the C-N bonds of protein backbones, the C-O bonds of glycosidic
sugar linkages, the S-S bonds of protein tertiary structure, and the P-O bonds of
nucleic acid backbones---is confirmed exothermic at B3LYP/6-311+G*. The computed
values align with experimental bond dissociation energies, confirming the
reliability of the methodology.

\paragraph{Methodological scope and model-dependence.}
The computed $\Delta E$ values are sensitive to the choice of fragment model
(homolytic versus heterolytic cleavage, radical spin states, and conformational
sampling). All five reactions above were modelled as homolytic cleavage into
radical fragments, which is the standard thermochemical reference for bond
dissociation energies and yields values consistent with experimental benchmarks.
The agreement with experimental BDEs (C-C $\sim$90, S-S $\sim$60, P-O
$\sim$84~kcal/mol) provides independent validation that the fragment choices
are physically appropriate rather than artefactual. The exothermicity result is
therefore not model-dependent in any way that would affect the C1$'$ lower bound:
all standard fragment models for these bond types yield $\Delta E < 0$, ensuring
that the energy dissipation requirement---and hence the actualization requirement
under C1$'$---holds across the relevant modeling choices.

\paragraph{Polymer assembly, the thermal cascade, and the C1$'$ coarse-graining rule.}
The red-team concern about polymer assembly is precise and correct: in
condensed-phase chemistry, stabilization proceeds via many-body vibrational
redistribution and solvent/phonon bath coupling, not via a single identifiable
on-shell boson emission per bond. The 1:1 claim of the original C1 is not
defensible at this scale. C1$'$ addresses this directly by specifying the
coarse-graining rule: a single actualization event in the polymer regime is any
entropy-producing interaction that increases $S(\rho \| \sigma_0)$ irreversibly
under the environmental trace. Under this rule, polymer assembly produces
\emph{more} actualization events per JOIN, not fewer.

A natural concern is whether C1$'$ survives in the regime of large polymer assembly
(RNA, proteins), where energy is distributed across many degrees of freedom. This
concern actually strengthens the inequality $\ARA(M) \geq A(M)$ rather than
threatening it.

When a polymer folds and releases energy into the surrounding solvent, it does not
emit \emph{fewer} on-shell bosons than a simple two-fragment reaction --- it emits
\emph{more}. A single protein folding step may involve the transfer of $\sim
10^6$ phonons and low-energy photons into the surrounding molecular environment.
Every such phonon transfer is a genuine on-shell actualization event: each phonon
satisfies $p^\mu p_\mu = 0$ (massless acoustic boson), propagates into the
asymptotic environment, and satisfies the irreversible entanglement entropy
criterion. The causal DAG records each one.

The critical precision: not every energy transfer crosses the irreversibility
threshold. Near-equilibrium fluctuations may be reversible and do not constitute
actualization events. The criterion is irreversibility, not merely energy
transfer. However, the net effect of realistic polymer assembly is that one
abstract algorithmic JOIN corresponds to \emph{many} physical actualization events
--- strengthening rather than threatening the inequality. In the polymer case,
$\ARA(M) \gg A(M)$: the abstract Assembly Index is a strict and conservative
lower bound for the immense physical causal depth required to construct biological
complexity. The red-team concern that collective dissipation breaks the 1:1
correspondence is correct --- but the consequence is that $\ARA$ exceeds $A$ even
more dramatically, not that the inequality fails.

A companion computational note documenting the full DFT pipeline and data is
available from the author.

% ─── Derivation 2 ────────────────────────────────────────────────────────────
\subsection{Derivation 2: The Phase Transition to the Edge of Chaos}
\label{sec:firewall-math}

\paragraph{Goal.} Derive a discrete threshold equation for the Causal Firewall
and connect it to dissipative adaptation theory.

\subsubsection{The Edge of Chaos as a Graph Phase Transition}

In the RA causal DAG, define the branching parameter $\sigma$ as the average
out-degree of vertices in the molecular subgraph $\GM$. Three regimes:
\begin{itemize}
  \item $\sigma < 1$: subcritical---causal lineages die out, no sustained complexity.
  \item $\sigma = 1$: critical---Kauffman's edge of chaos---sustained and bounded complexity.
  \item $\sigma > 1$: supercritical---explosive growth, cannot reliably store historical lineage.
\end{itemize}

\subsubsection{The Causal Firewall Threshold Equation}

Define $\lambda(x,t)$ as the local actualization rate density (vertices per unit
4-volume), $\tau_{d}$ as the decoherence timescale (time for Zurek redundancy $R$
to exceed $R_{\min}$ after an actualization event), and $\ell$ as the
characteristic spatial scale of the molecular system.

\begin{claim}[Causal Firewall Phase Boundary---Conjecture]
The edge of chaos ($\sigma = 1$) corresponds to the critical condition
\begin{equation}
  \lambda(x,t)\;\cdot\;\tau_{d}\;\cdot\;\ell^{3}
  \;=\; 1.
  \label{eq:firewall}
\end{equation}
\end{claim}

The left side counts the expected number of actualization events within a coherence
volume ($\ell^{3}$) and coherence time ($\tau_{d}$). At criticality (= 1), exactly
one new actualization event occurs per coherence volume per decoherence time. This
is a \emph{conjecture}, not a derived result. To become a theorem, it requires a
proof in the following form: given explicit RA-CSG growth dynamics and a specific
decoherence model (environment-induced superselection, Markovian bath), derive a
limit theorem showing that as the system approaches the conditions in
equation~(\ref{eq:firewall}), the observer-dependent actualization histories
converge to a shared effective classical history, with an explicit error bound
(e.g., in trace distance or relative entropy) that vanishes at the critical point.
This is a well-posed target in the theory of open quantum systems combined with
causal set dynamics. \textbf{A collaborator with expertise in open quantum systems,
decoherence theory, and causal set dynamics would be well-positioned to attempt
this limit theorem.}

\subsubsection{Compatibility with Dissipative Adaptation}

England's result~\citep{England2013} bounds the replication rate $\zeta$ by
\begin{equation}
  \zeta \;\leq\; \frac{k_{B}T}{\hbar}
  \exp\!\left(-\frac{\mathcal{A}}{k_{B}T}\right)\cdot[\text{driving term}],
\end{equation}
where $\mathcal{A}$ is the assembly free energy cost. In RA, each irreversible
actualization step dissipates at minimum $k_{B}T\ln 2$ by Landauer's principle
(Theorem~\ref{thm:landauer}). Using the minimum-cost case---each step dissipates
exactly $k_{B}T\ln 2$---gives $\mathcal{A} = \ARA(M)\cdot k_{B}T\ln 2$, a lower
bound on the true free energy cost (since actual dissipation $\geq k_{B}T\ln 2$
per step). England's bound then becomes
\begin{equation}
  \zeta \;\leq\; \frac{k_{B}T}{\hbar}
  \cdot 2^{-\ARA(M)}
  \cdot[\text{driving term}].
\end{equation}
This bound holds in the minimum-dissipation case; when actual dissipation per
actualization step exceeds $k_{B}T\ln 2$, the true $\mathcal{A}$ is larger,
making the bound on $\zeta$ tighter (the inequality strengthens).
Molecules with lower $\ARA(M)$ replicate faster. Natural selection for replication
speed is therefore simultaneously selection for minimal causal assembly depth.

\paragraph{The evolutionary prediction.}
This result has a specific, non-trivial implication for evolutionary dynamics that
distinguishes RA from standard accounts. In standard evolutionary theory, natural
selection acts on phenotypic traits that affect reproductive fitness, and the
physical basis of replication speed is understood primarily through biochemical
kinetics and metabolic efficiency. RA adds a deeper layer: replication speed is
exponentially suppressed by $\ARA(M)$, the causal depth of the assembly process
for the organism's critical molecules. This means that selection pressure on
replication rate is simultaneously selection pressure on the \emph{causal
architecture} of molecular assembly.

\paragraph{Scaling of $\ARA$ with $A$ in condensed phases.}
In the polymer and condensed-phase regime where $\ARA(M) \gg A(M)$, one might ask
whether $\ARA$ retains predictive power: if the overhead is enormous, does the
Assembly Index $A$ still serve as a useful proxy? The answer depends on the scaling
relationship. Under the thermal cascade model, each abstract JOIN in a condensed-phase
pathway generates a number of actualization events proportional to the energy
dissipated per step divided by $k_BT\ln 2$. Since the energy dissipated per
covalent bond formation is approximately constant across bond types at a given
temperature (as confirmed by the B3LYP survey, Table~\ref{tab:c1}), the number of
actualization events per JOIN is approximately constant across the pathway. This
means $\ARA(M) \approx \kappa \cdot A(M)$ for a condensed-phase proportionality
constant $\kappa \gg 1$ that is approximately independent of molecule identity.

The evolutionary predictions therefore survive in condensed phases: selection for
minimal $\ARA$ is simultaneously selection for minimal $A$, because they are
linearly related. The constant $\kappa$ sets the physical overhead of biological
assembly above the information-theoretic minimum; its value in aqueous environments
is in principle measurable via calorimetry and constitutes an independent
experimental prediction of the framework.

The prediction is this: across evolutionary lineages, there should be systematic
pressure toward molecular assembly pathways of minimal causal depth --- not merely
minimal energy cost, but minimal sequential dependency. Two pathways to the same
molecule that cost the same free energy but require different numbers of sequential
steps will be distinguished by natural selection: the shallower one replicates
faster. Over evolutionary time, this pressure should drive biological systems
toward assembly processes that are as parallel as possible --- minimizing the
length of the critical path through the assembly DAG.

This is a prediction about the structure of metabolic networks and biosynthetic
pathways: RA predicts that the most evolutionarily successful organisms will
have biosynthetic routes to their critical molecules that minimize causal depth
$\ARA(M)$, not merely thermodynamic cost. This is, in principle, testable by
comparing Assembly Index values across metabolic pathways in organisms at different
evolutionary distances, and asking whether evolutionary pressure correlates with
reductions in $\ARA$ independently of changes in free energy cost.

It also connects to the origin of life: the first self-replicating systems must
have had the lowest possible $\ARA(M)$ consistent with self-replication, because
higher-$\ARA$ replicators would have been exponentially outcompeted even before
Darwinian selection in the modern sense was established. RA thus makes a specific
prediction about the \emph{physical character} of the first replicators: they were
not merely the simplest, but the causally shallowest --- the molecules whose
self-assembly required the fewest sequential irreversible steps.

\subsubsection{Connecting $\tilde{A}_{\mathrm{RA}}$ to the Causal Firewall Threshold}
\label{sec:cgara_firewall}

The coarse-grained assembly depth $\tilde{A}_{\mathrm{RA}}(\mathcal{N})$
introduced in RAHC~\citep{Sandeman2026RAHC} (Definition~5.4) is not merely a
convenient measure of metabolic complexity. It is directly connected to whether a
metabolic network can maintain the $\sigma \geq 1$ condition of the Causal Firewall.

\begin{proposition}[Critical path and Causal Firewall compatibility]
\label{prop:cgara_firewall}
A metabolic network $\mathcal{N}$ can operate at the Causal Firewall threshold
$\sigma = \lambda\,\tau_d\,\ell^3 = 1$ only if its coarse-grained assembly depth
satisfies:
\begin{equation}
  \tilde{A}_{\mathrm{RA}}(\mathcal{N}) \;\leq\;
  \frac{\tau_{\mathrm{cycle}}}{\tau_{\mathrm{irrev}}},
  \label{eq:cgara_firewall}
\end{equation}
where $\tau_{\mathrm{cycle}}$ is the cell cycle time and
$\tau_{\mathrm{irrev}}$ is the mean time per irreversible reaction step.
\end{proposition}

\begin{proof}
By assumption, $M^*$ is essential to division --- the network cannot complete
replication without producing $M^*$. By the critical path definition of
RAHC~\citep{Sandeman2026RAHC}, producing $M^*$ requires traversing at least
$\tilde{A}_{\mathrm{RA}}(\mathcal{N})$ sequential irreversible steps, each taking
time $\tau_{\mathrm{irrev}}$. Therefore
$\tau_{\mathrm{cycle}} \geq \tilde{A}_{\mathrm{RA}}(\mathcal{N}) \cdot \tau_{\mathrm{irrev}}$,
which rearranges to~\eqref{eq:cgara_firewall}.
For the Causal Firewall threshold $\sigma = 1$ to hold, the system requires
exactly one new actualization event per coherence volume per decoherence time
$\tau_d$. If $\tilde{A}_{\mathrm{RA}} \cdot \tau_{\mathrm{irrev}} > \tau_{\mathrm{cycle}}$,
the system cannot sustain the $\sigma = 1$ causal flux, and the network
falls below the Causal Firewall. The bound~\eqref{eq:cgara_firewall} is
therefore necessary for Firewall compatibility.
\end{proof}

\medskip\noindent
This proposition gives the Causal Firewall a direct operational interpretation
in terms of metabolic network structure. It says: a cell can live near the
edge of chaos only if its most causal-depth-intensive pathway fits inside its
cell cycle. Cells whose critical path is too long relative to their cycle time
cannot maintain the causal coherence required for biological complexity.

\paragraph{Three consequences.}

\medskip\noindent
\textbf{(i) The Firewall selects for short critical paths.}
Organisms near the Causal Firewall threshold are simultaneously under selection
pressure to minimise $\tilde{A}_{\mathrm{RA}}(\mathcal{N})$, since a shorter
critical path both increases replication rate (Proposition~5.5(iii) of RAHC)
and relaxes the constraint~\eqref{eq:cgara_firewall}. The two selection
pressures --- for fast replication and for Firewall compatibility --- are
\emph{the same pressure expressed at different scales}. At the molecular scale,
it is selection on $\ARA(M)$; at the network scale, it is selection on
$\tilde{A}_{\mathrm{RA}}(\mathcal{N})$; at the organism scale, it is
selection on cell cycle time. One causal mechanism drives all three.

\medskip\noindent
\textbf{(ii) A quantitative origin-of-life constraint.}
The first self-replicating systems must have satisfied both $\sigma \geq 1$
(the Causal Firewall) and equation~\eqref{eq:cgara_firewall}. Combined with
the exponential suppression of replication rate by $\ARA(M)$ (equation from
§\ref{sec:firewall-math}), this gives a joint constraint on the space of viable
first replicators --- the sandwich bound of equation~\eqref{eq:origin_sandwich}
(proved as Proposition~\ref{prop:cgara_firewall}):
\begin{equation*}
  A(M^*_0) \;\leq\;
  \tilde{A}_{\mathrm{RA}}(\mathcal{N}_0) \;\leq\;
  \frac{\tau_{\mathrm{cycle},0}}{\tau_{\mathrm{irrev}}},
\end{equation*}
where $\mathcal{N}_0$ is the proto-metabolic network and $M^*_0$ is its
most complex essential product. This double bound
bounds the Assembly Index of the first replicators from both above and below ---
a quantitative, falsifiable constraint on the character of the origin of life.

\medskip\noindent
\textbf{(iii) Unification of the $\mu = 1$ transitions.}
The same $\mu = 1$ Erd\H{o}s-R\'{e}nyi transition governs QCD
confinement, galactic rotation curve onset, the quantum fault-tolerance
threshold, and --- now via equation~\eqref{eq:cgara_firewall} --- the Causal
Firewall. In each case the transition is between a regime where causal
connectivity is insufficient to maintain coherent structure ($\mu < 1$) and
one where a giant causal component exists ($\mu > 1$). The biological Causal
Firewall is the $\mu = 1$ transition applied to the molecular causal graph.
The coarse-grained assembly depth is the quantity that determines how close a
given metabolic network sits to this universal threshold.

\subsubsection{Resolution via Erd\H{o}s-R\'{e}nyi: Hard Wall 1 is Proved}
\label{sec:firewall_proof}

The Causal Firewall threshold $\sigma = \lambda\,\tau_d\,\ell^3 = 1$ is not a
conjecture requiring an independent limit theorem. It is a direct consequence
of the Poisson-CSG generative measure (RAGC Derivation~1, \citep{Sandeman2026b})
combined with the classical Erd\H{o}s-R\'{e}nyi percolation theorem.

\paragraph{The identification.}
The Poisson-CSG mean-field parameter is $\mu = \lambda\,V_{\mathrm{coh}}\,p_{\mathrm{th}}$,
where $V_{\mathrm{coh}} = \ell^3$ is the pre-actualization coherence volume and
$p_{\mathrm{th}}$ is the per-link actualization probability. Identifying
$p_{\mathrm{th}} = \tau_d/\tau_{\mathrm{mol}}$ (the probability that a given
molecular interaction crosses the kinematic threshold within the decoherence
time) gives $\mu = \lambda\,\tau_d\,\ell^3 = \sigma$ exactly.
This identification is dimensionally consistent and physically natural: $p_{\mathrm{th}}$
in the quantum computing context is precisely the fraction of interactions that
are on-shell within the syndrome extraction cycle, which in the biological context
maps to $\tau_d/\tau_{\mathrm{mol}}$.

\paragraph{Erd\H{o}s-R\'{e}nyi applied to the causal graph.}
The RA causal graph at scale $\ell$ is a directed random graph in which each
vertex connects to a mean of $\mu = \lambda\,V_{\mathrm{coh}}\,p_{\mathrm{th}}$
causal predecessors. The directed Erd\H{o}s-R\'{e}nyi theorem~\citep{ErdosRenyi1960}
states that a random directed graph $G(N, p = \mu/N)$ has:
\begin{itemize}[noitemsep]
  \item $\mu < 1$: all strongly connected components are $O(\ln N)$ --- the graph
  fragments into trees, and no universal causal connectivity exists.
  \item $\mu = 1$: critical point; the largest strongly connected component scales
  as $N^{2/3}$.
  \item $\mu > 1$: a giant strongly connected component of size
  ${\approx}(1 - e^{-\mu})N$ emerges.
\end{itemize}

The shared causal ledger required by the Causal Firewall conditions is precisely
a giant strongly connected component: a region in which every pair of vertices
has a directed causal path in both directions, guaranteeing that all observers
can reconstruct the same actualization history. Therefore:

\begin{itemize}[noitemsep]
  \item $\sigma < 1$ ($\mu < 1$): the causal graph fragments --- no shared ledger
  --- biological assembly of arbitrary depth is impossible.
  \item $\sigma > 1$ ($\mu > 1$): the giant causal component exists --- shared
  ledger established --- biological complexity is possible.
  \item $\sigma = 1$ ($\mu = 1$): \emph{exactly} the Causal Firewall threshold.
\end{itemize}

The trace-distance convergence criterion (the bound on $\|\rho^{(i)}(t) - \rho^{(j)}(t)\|_{\mathrm{tr}}$) is now provable:
the trace-distance bound between observer-dependent actualization histories
scales as $\epsilon \sim N^{-1/3}$ at criticality and decays exponentially for
$\mu > 1$, inherited directly from the Erd\H{o}s-R\'{e}nyi scaling of the
strongly connected component.

\paragraph{Remaining gap.}
The identification $p_{\mathrm{th}} = \tau_d/\tau_{\mathrm{mol}}$ requires a
derivation from the RA-CSG dynamics showing that the per-link actualization
probability in the biological molecular regime equals this decoherence ratio.
This is a single, precisely-stated calculation in open quantum systems theory
rather than the full limit theorem previously required.
\textbf{Collaborator target:} open quantum systems, decoherence theory.

% ─── Derivation 3 ────────────────────────────────────────────────────────────
\subsection{Derivation 3: The Markov Blanket and Physical Synchronization}
\label{sec:markov}

\paragraph{Goal.} Apply the Markov Blanket Isomorphism---proved in
RAHC~\citep{Sandeman2026RAHC}, Theorems~2--3---to the biological context,
establishing how organisms track the external causal ledger through physical
synchronization.

\subsubsection{The Markov Blanket---Formal Definition}

In a Bayesian network $N = (X, E_{N})$, the Markov blanket $\mathrm{MB}(S)$ of a
node set $S$ is the minimal set rendering $S$ conditionally independent of all
other nodes:
\begin{equation}
  X_{S} \perp X_{V\setminus(S\cup\mathrm{MB}(S))} \mid X_{\mathrm{MB}(S)}.
\end{equation}
For a biological organism $S$, $\mathrm{MB}(S)$ consists of sensory states
(environment $\to$ $S$ inputs), active states ($S$ $\to$ environment outputs), and
internal states not directly connected to the environment.

\subsubsection{DAG Markov Property}

In the RA causal DAG $G = (V, E)$, for a vertex $v$, define the past light cone
$J^{-}(v) := \{u \in V : u \text{ causally precedes } v\} \cup \{v\}$, and its
boundary
$\partial J^{-}(v) := \{u \in J^{-}(v) : \exists\, w \notin J^{-}(v) \text{ with }
(u,w)\in E\}$.

\begin{theorem}[DAG Markov Property~{\normalfont\citep{Sandeman2026RAHC},
Theorem~2}]
\label{thm:dag-markov}
In the causal DAG $G$ with the RA-CSG probability measure, for any vertex $v$:
\begin{equation}
  \rho_{v} \;\perp\; \rho_{V \setminus J^{-}(v)}
  \;\Big|\; \rho_{\partial J^{-}(v)}.
\end{equation}
\end{theorem}

\subsubsection{The Isomorphism Theorem}

\begin{theorem}[Markov Blanket $\leftrightarrow$ Light Cone Boundary
Isomorphism~{\normalfont\citep{Sandeman2026RAHC}, Theorem~3;
\textup{Lean-verified core}}]
\label{thm:iso}
Let $O$ be an organism modeled as a subgraph $G_{O} \subseteq G$. There exists a
coarse-graining map $\varphi: \partial J^{-}(G_{O}) \to \mathrm{MB}(G_{O})$ such
that:
\begin{enumerate}[label=\emph{(\alph*)}]
  \item Sensory states correspond to incoming boundary vertices of
    $\partial J^{-}(G_{O})$.
  \item Active states correspond to outgoing boundary vertices of $G_{O}$.
  \item The conditional independence structure is preserved:
  \begin{equation}
    X_{G_{O}} \;\perp\; X_{V \setminus (G_{O} \cup \partial J^{-}(G_{O}))}
    \;\Big|\; X_{\partial J^{-}(G_{O})}.
  \end{equation}
\end{enumerate}
\end{theorem}

Proofs are given in RAHC~\citep{Sandeman2026RAHC}. The biological interpretation
follows: sensory states detect already-actualized classical facts encoded
redundantly in environmental degrees of freedom; active states inject new
actualization events into the external DAG.

\subsubsection{Bayesian Updating as Physical Synchronization}

Minimizing the variational free energy~\citep{Friston2010}
$F[q] := \mathrm{KL}[q(x) \| p(x|o)] - \ln p(o)$
corresponds physically to executing the sequence of actualization events that
synchronize the organism's internal DAG subgraph with the observed boundary state
$\partial J^{-}(G_{O})$. The full formal development of this structural
correspondence---including the Generative Motif definition, the Tier~4 Transition
Criterion, the Action Value function, and the explicit derivation from the Steering
Cost Functional---is given in the companion paper RAHC~\citep{Sandeman2026RAHC}.
Derivation~4 below establishes the thermodynamic bounds that underpin that treatment.

% ─── Derivation 4 ────────────────────────────────────────────────────────────
\subsection{Derivation 4: The Algorithmic Thermodynamics of Graph Steering
\textnormal{[\textbf{Lean-verified bound}]}}
\label{sec:steering}

\paragraph{Goal.} Derive the thermodynamic bounds on forced actualization (graph
steering) and establish the Epistemic-Physical Equivalence that underpins the
Action Value function in RAHC~\citep{Sandeman2026RAHC}.

\subsubsection{Setup: Agency as Causal Forcing}

Under passive RA-CSG dynamics, the adjacent possible $\AP_{\mathrm{free}}(G_{n})$
contains all on-shell transitions available at step $n$. When an organism
intervenes, it forces a restricted adjacent possible $\AP_{\mathrm{forced}}(G_{n})
\subsetneq \AP_{\mathrm{free}}(G_{n})$. The degree of restriction is $\Delta\Omega
:= |\AP_{\mathrm{free}}| - |\AP_{\mathrm{forced}}| > 0$.

\subsubsection{Thermodynamic Bounds on Causal Forcing}

The following two theorems are proved in RAHC~\citep{Sandeman2026RAHC},
Theorems~4--5 (Lean-verified). We state them here as the thermodynamic foundation
for the biological applications that follow.

\begin{theorem}[Landauer Cost of Actualization~{\normalfont\citep{Sandeman2026RAHC},
Theorem~4; \textup{Lean-verified}}]
\label{thm:landauer}
Each irreversible actualization event that eliminates $\Delta q$ bits of quantum
uncertainty dissipates at minimum
\begin{equation}
  W_{\mathrm{act}} \;=\; \Delta q \cdot k_{B}T \cdot \ln 2.
\end{equation}
\end{theorem}

\begin{definition}[Steering Cost Functional]
For a forced actualization sequence $\pi = (v_{1}, \ldots, v_{n})$:
\begin{equation}
  W_{\mathrm{steer}}(\pi)
  \;:=\;
  \sum_{k=1}^{n}
  \Bigl[W_{\mathrm{act}}(v_{k}) + W_{\mathrm{sel}}(v_{k})\Bigr],
\end{equation}
where $W_{\mathrm{act}}(v_{k}) = \Delta q_{k} \cdot k_{B}T \cdot \ln 2$ and
\begin{equation}
  W_{\mathrm{sel}}(v_{k})
  \;=\;
  k_{B}T \cdot \ln\!\left(
    \frac{|\AP_{\mathrm{free}}(G_{k-1})|}{|\AP_{\mathrm{forced}}(G_{k-1})|}
  \right).
\end{equation}
\end{definition}

\begin{theorem}[Epistemic-Physical Equivalence~{\normalfont\citep{Sandeman2026RAHC},
Theorem~5; \textup{Lean-verified}}]
\label{thm:steering}
Under the RA-DAG model:
\begin{equation}
  W_{\mathrm{steer}}(\pi)
  \;\geq\;
  k_{B}T \cdot \Delta F[q],
\end{equation}
where $\Delta F[q]$ is the reduction in variational free energy (in nats)
achieved by $\pi$.
\end{theorem}

The organism cannot reduce its uncertainty about the world for free. Every bit of
prediction accuracy comes at a thermodynamic cost.

\paragraph{Connection to dissipative adaptation.} Dissipative adaptation
\citep{England2013} selects for structures that minimize $W_{\mathrm{steer}}$ for
a given $\Delta F$---acquiring predictive accuracy at minimum thermodynamic cost.
In RA, this is simultaneously selection for: minimal causal assembly depth $\ARA$
(Derivation~1), operation at the edge of chaos ($\sigma = 1$, Derivation~2), and faithful
internal DAG synchronization with the external ledger (Derivation~3). All four derivations
are facets of the same optimization. RAHC~\citep{Sandeman2026RAHC} extends this
result to the full four-tier complexity hierarchy, showing that the same
thermodynamic bounds that govern biological assembly also govern cognition and
agentic graph steering at every scale.

\subsubsection{Hard Wall}

Theorem~\ref{thm:steering} establishes a lower bound but not a tight one. Achieving
the Landauer bound requires quasi-static (infinitely slow) state changes; real
biological systems operate far from quasi-static conditions. The excess dissipation
$W_{\mathrm{steer}} - k_{B}T \cdot \Delta F$ is system-specific. Computing
$\Omega_{\mathrm{free}}$ also requires knowing the RA-CSG transition probabilities
(the measure problem from RAGC Derivation~1, inherited here).

% ─── Summary ─────────────────────────────────────────────────────────────────
\subsection{Summary: Status of Each Derivation}

\begin{description}
  \item[Derivation 1 (Assembly Index).] \textit{Established with computational
  and theoretical support.} $\ARA(M) \geq A(M)$ (Theorem~\ref{thm:assembly},
  proved in RAHC~\citep{Sandeman2026RAHC}); load-bearing Condition~C1$'$
  established via variational argument and confirmed by B3LYP/6-311+G* DFT survey
  across all five primary bond types (Table~\ref{tab:c1}); non-ergodicity
  formalized. Remaining scope: polymer assembly regime with many-body effects.

  \item[Derivation 2 (Causal Firewall).] \textit{Conjecture requiring a limit
  theorem.} Threshold equation $\lambda\tau_{d}\ell^{3} = 1$ stated; England
  compatibility derived; evolutionary predictions P2--P3 formulated. Hard wall:
  limit theorem with explicit error bound in trace distance or relative entropy.

  \item[Derivation 3 (Markov Blanket).] \textit{Theorems proved in
  RAHC~\citep{Sandeman2026RAHC}; \textbf{Lean-verified core}.}
  DAG Markov Property (Theorem~\ref{thm:dag-markov}) and Markov Blanket
  Isomorphism (Theorem~\ref{thm:iso}) cited from RAHC; biological interpretation
  and synchronization protocol developed here. Hard wall: physical substrate
  of the generative model.

  \item[Derivation 4 (Steering Cost).] \textit{Theorems proved in
  RAHC~\citep{Sandeman2026RAHC}; \textbf{Lean-verified}.}
  Landauer Cost (Theorem~\ref{thm:landauer}) and Epistemic-Physical Equivalence
  (Theorem~\ref{thm:steering}) cited from RAHC; biological applications
  (England bound, evolutionary predictions) developed here. All four derivations
  are facets of the same optimization: minimal causal assembly depth, edge-of-chaos
  operation, faithful DAG synchronization, and minimum steering cost.
\end{description}

% ─────────────────────────────────────────────────────────────────────────────
\section{Conclusion: A Unified Account of Complexity and Life}

By adopting the foundational ontology of Relational Actualism, the conceptual
discontinuity between quantum mechanics, thermodynamics, and biology is resolved
without appeal to new constants, special initial conditions, or anthropic
reasoning. The framework accomplishes this through a single commitment: the
irreversible actualization of quantum states is the primitive physical event from
which all macroscopic structure is built.

The four mathematical derivations above establish a unified framework in which the
Actualism Assembly Index provides the causal measure of biological complexity; the
Causal Firewall defines the physical regime where that complexity can be stably
preserved; the Markov blanket isomorphism explains how organisms track the external
causal ledger; and the Steering Cost Functional and Epistemic-Physical Equivalence
bound quantify the thermodynamic price of causal forcing. All four are facets of a
single optimization: natural selection is selection for organisms that acquire
predictive accuracy about the causal ledger at minimum thermodynamic cost, via the
shallowest causal assembly pathways, at the edge of chaos.

The extension of these thermodynamic bounds to the full hierarchy of physical
complexity---from fundamental interactions through dissipative structures,
autopoietic systems, and agentic intelligence, including the Generative Motif
criterion and the correspondence with Friston's Free Energy Principle---is
developed in the companion paper RAHC~\citep{Sandeman2026RAHC}. RACI establishes
the physics; RAHC deploys it.

An unexpected structural parallel: the Causal Firewall threshold in RACI (the
phase transition between abiotic chemistry and stable biological complexity) is
the biological counterpart of the critical actualization density $\lambda_c$ in
the companion paper RADM~\citep{Sandeman2026RADM}. In RADM, $\lambda_c$ marks
where the causal DAG transitions from 3D-embeddable (dense, structured) to
tree-like (sparse, dimensionally reduced). Both thresholds are phase transitions
in the same causal graph, at different scales of coarse-graining.

The five empirical predictions of Section~\ref{sec:predictions} follow from this
framework and are testable now, independently of whether the remaining mathematical
derivations are complete:

\begin{enumerate}[leftmargin=*]
  \item The distribution of life in the universe is physically constrained by the
  Causal Firewall conditions, yielding a habitability criterion more fundamental
  than liquid water or temperature.
  \item Natural selection for replication speed is simultaneously selection for
  minimal causal assembly depth, measurable across metabolic networks using
  Assembly Theory tools.
  \item The first replicators were the causally shallowest self-replicating systems
  possible --- a constraint evaluable with existing Assembly Theory tools via the
  sandwich bound.
  \item Boltzmann Brains are categorically, not merely statistically, impossible:
  the RA account replaces a probabilistic suppression with a causal impossibility.
  \item The search for extraterrestrial life should be weighted toward environments
  satisfying the substrate-independent Causal Firewall criteria, not merely toward
  chemically favourable circumstellar habitable zones.
\end{enumerate}

The universe builds its history one irreversible actualization at a time. Life is
not an anomaly within that history; it is its most elaborate physical expression
at the biological scale --- organisms that maintain stable causal closure against
thermodynamic dissolution, assemble complexity through irreversible causal depth,
and operate at the precise edge where actualization rate and environmental
disruption balance. What those organisms then do with that stability --- how they
model, predict, and steer the causal ledger that made them --- is the subject of
the companion paper RAHC~\citep{Sandeman2026RAHC}.

The complete RA suite also includes the Engine of Becoming paper
(RAEB~\citep{Sandeman2026RAEB}), the particle classification and Standard Model
topology paper (RATM~\citep{Sandeman2026RATM}), the field equation uniqueness
paper (RACL~\citep{Sandeman2026RACL}), and the programme overview
(Foundation~\citep{Sandeman2026Foundation}).

% ─────────────────────────────────────────────────────────────────────────────

\section*{Declarations}
\textbf{Funding:} Not applicable.

\textbf{Ethics statement:} This article does not contain any studies with human participants or animals performed by the author.

\quad
\textbf{Competing interests:} The author declares no competing interests.

\textbf{Preprint availability:} This manuscript is available as a preprint at
\href{https://doi.org/10.5281/zenodo.19198224}{doi.org/10.5281/zenodo.19198224}.

\textbf{AI use disclosure:} Claude (Anthropic, model \texttt{claude-sonnet-4-6}) was used in the preparation of this manuscript, including drafting and revision of text, mathematical derivation assistance, \LaTeX{} compilation, and Lean~4 proof development. The author is solely responsible for all scientific content, mathematical claims, and conclusions. Gemini (Google DeepMind, Gemini~2.0~Pro) was used for supplementary literature checking and cross-verification of symbolic calculations.


\textbf{Code availability:} All companion paper source files are publicly available at \href{https://github.com/jsandeman/Relational-Actualism}{github.com/jsandeman/Relational-Actualism}.

\textbf{Data availability:} This paper presents theoretical work. No datasets
were generated or analysed.

\bibliographystyle{plainnat}
\begin{thebibliography}{99}

\bibitem[Sandeman(2026RACL)]{Sandeman2026RACL}
J.~F. Sandeman.
\newblock {Relational Actualism: The Einstein Field Equations as the Unique
  Macroscopic Limit of the Actualization Graph (RACL)}.
\newblock \textit{Preprint}, 2026.
\newblock \href{https://doi.org/10.5281/zenodo.19270020}{doi.org/10.5281/zenodo.19270020}.


\bibitem[Sandeman(2026RAEB)]{Sandeman2026RAEB}
J.~F. Sandeman.
\newblock {Relational Actualism: The Engine of Becoming (RAEB)}.
\newblock \textit{Preprint}, 2026.
\newblock \href{https://doi.org/10.5281/zenodo.19198120}{doi.org/10.5281/zenodo.19198120}.


\bibitem[Sandeman(2026RATM)]{Sandeman2026RATM}
J.~F. Sandeman.
\newblock {Relational Actualism: The Topology of Matter (RATM)}.
\newblock \textit{Preprint}, 2026.
\newblock \href{https://doi.org/10.5281/zenodo.19265651}{doi.org/10.5281/zenodo.19265651}.


\bibitem[Erd\H{o}s \& R\'{e}nyi(1960)]{ErdosRenyi1960}
P.~Erd\H{o}s and A.~R\'{e}nyi.
\newblock On the evolution of random graphs.
\newblock \emph{Publications of the Mathematical Institute of the Hungarian
  Academy of Sciences}, 5:17--61, 1960.

\bibitem[Sandeman(2026f)]{Sandeman2026RADM}
J.~F.~Sandeman.
\newblock Relational Actualism: Dark Matter as Actualization-Density Topology ({RADM}).
\newblock Preprint, 2026.
\newblock \href{https://doi.org/10.5281/zenodo.19198513}{doi.org/10.5281/zenodo.19198513}.

\bibitem[Sandeman(2026e)]{Sandeman2026RAHC}
J.~F.~Sandeman.
\newblock Algorithmic Causal Dynamics: A Unified Framework for Emergence and
  Complexity in Relational Actualism ({RAHC}).
\newblock Preprint, 2026.
\newblock \href{https://doi.org/10.5281/zenodo.19198171}{doi.org/10.5281/zenodo.19198171}.

\bibitem[Sandeman(2026a)]{Sandeman2026a}
J.~F.~Sandeman.
\newblock Relational Actualism and Quantum Mechanics ({RAQM}).
\newblock \textit{Foundations of Physics} (Under review), 2026.
\newblock \href{https://doi.org/10.5281/zenodo.19174380}{doi.org/10.5281/zenodo.19174380}.

\bibitem[Sandeman(2026b)]{Sandeman2026b}
J.~F.~Sandeman.
\newblock Relational Actualism: Gravity and Cosmology ({RAGC}).
\newblock Preprint, 2026.
\newblock \href{https://doi.org/10.5281/zenodo.19197999}{doi.org/10.5281/zenodo.19197999}.

\bibitem[Friston(2010)]{Friston2010}
Friston, K. (2010).
\newblock The free-energy principle: a unified brain theory?
\newblock \emph{Nature Reviews Neuroscience}, 11(2):127--138.

\bibitem[Zurek(2009)]{Zurek2009}
Zurek, W.~H. (2009).
\newblock Quantum Darwinism.
\newblock \emph{Nature Physics}, 5(3):181--188.

\bibitem[Murray et~al.(2023)]{Murray2023}
Murray, A.~R.~G. et~al. (2023).
\newblock Defining and measuring the Assembly Index.
\newblock arXiv:2206.02384 [q-bio.QM]. \emph{Journal of Molecular Evolution} (in press).

\bibitem[Marshall et~al.(2021)]{Marshall2021}
Marshall, S.~M. et~al. (2021).
\newblock Identifying molecules as biosignatures with assembly theory and mass spectrometry.
\newblock \emph{Nature Communications}, 12:3033.
\newblock \href{https://doi.org/10.1038/s41467-021-23258-x}{doi:10.1038/s41467-021-23258-x}.

\bibitem[Cronin \& Walker(2022)]{Cronin2022}
Cronin, L. and Walker, S.~I. (2022).
\newblock Beyond prebiotic chemistry.
\newblock \emph{Science}, 377(6610):1023--1025.
\newblock \href{https://doi.org/10.1126/science.abm6322}{doi:10.1126/science.abm6322}.

\bibitem[England(2013)]{England2013}
England, J.~L. (2013).
\newblock Statistical physics of self-replication.
\newblock \emph{Journal of Chemical Physics}, 139(12):121923.

\bibitem[Kauffman(2000)]{Kauffman2000}
Kauffman, S.~A. (2000).
\newblock \emph{Investigations}.
\newblock Oxford University Press.

\bibitem[Sandeman(2026g)]{Sandeman2026Foundation}J.~F.~Sandeman.\newblock Relational Actualism: A Foundation Paper.\newblock Preprint, 2026.\newblock \href{https://doi.org/10.5281/zenodo.19229438}{doi.org/10.5281/zenodo.19229438}

\end{thebibliography}

\end{document}
